\documentclass[twocolumn]{iopjournal}

\usepackage{amsmath,amssymb,amsfonts}

\fancyhfoffset[L]{0mm}
\setlength{\marginparwidth}{0pt}
\setlength{\marginparsep}{0pt}

% ==========================================
% ZENODO PREPRINT HEADERS (ACTIVE)
% ==========================================
\renewcommand{\articletype}[1]{%
  {\vspace*{-8mm}\noindent \Large \sf Preprint}\par
  \vspace*{6mm}\noindent{\scriptsize \sf{\bfseries \MakeUppercase{#1}}}}

\fancyhead[L]{}
\fancyhead[R]{}

\hypersetup{pdftitle={The Stellated Octahedron, the Cube and the Hadamard Matrix}, pdfauthor={Stephen Nelson}, pdfsubject={Gravity Field Dynamics (GFD)}, pdfkeywords={Hadamard matrix, Walsh functions, cube, stellated octahedron, tetrad, proof}}

\begin{document}

\articletype{Supplementary Material}

\title{The Stellated Octahedron, the Cube and the Hadamard Matrix}

\author{Stephen Nelson}

\affil{Field Dynamics Research}

\ead{stephen@fielddynamics.org}

\begin{abstract}
A cube simplex is a set of $D + 1$ vertices of the cube
$\{-1,+1\}^D$ that are pairwise equidistant. We show that the
coordinate matrix $M$ of a cube simplex, bordered by a row of ones,
is a Hadamard matrix, and conversely; that this is the forced Gram
matrix of the tetrad frame, $M^TM = (D+1)I - J$; and that the
Hadamard rows encode closure and an isotropic second moment. A cube
simplex therefore exists only if $D + 1$ is $2$ or divisible by $4$.
A cube simplex and its parity image are disjoint, and together they
fill the cube only for $D = 3$, where they are the two parity
classes $x_1x_2x_3 = \pm1$ and their compound is the stellated
octahedron. This is the content of the equation $2(D+1) = 2^D$. In a
stated ordering, the $8 \times 8$ Walsh matrix of the cube is exactly
$H_2 \otimes \hat H$, with $H_2$ acting as spatial parity and
$\hat H$ the bordered frame. Only for $K = 2$ and $K = 4$ do the rows
of $M$ exhaust the ways of splitting the $K$ points into equal halves.
\end{abstract}

\keywords{Hadamard matrix, Walsh functions, cube, stellated octahedron, tetrad, proof}

% ----------------------------------------------------------------------

\section*{Setting}

Let $Q_D = \{-1,+1\}^D$, $D \ge 1$, and $K = D + 1$. A \emph{cube
simplex} is a set of $K$ distinct points of $Q_D$ that are pairwise
equidistant. Write them as the columns $v_0, \dots, v_D$ of a
$D \times K$ matrix $M$ with entries $\pm1$, let
$\mathbf{1} \in \mathbb{R}^K$ be the all-ones vector, and let
\begin{equation}
  \hat H = \begin{pmatrix} \mathbf{1}^T \\ M \end{pmatrix},
  \label{eq:Hhat}
\end{equation}
a $K \times K$ matrix of signs. A Hadamard matrix of order $n$ is an
$n \times n$ matrix $H$ of signs with $H^TH = nI$; for a square
matrix this is equivalent to $HH^T = nI$.

% ----------------------------------------------------------------------

\section*{The Frame Is a Hadamard Matrix}

\paragraph{Theorem 1.} \textit{For a $D \times K$ sign matrix $M$
the following are equivalent:}
\begin{itemize}
\item[(a)] \textit{the columns of $M$ form a cube simplex;}
\item[(b)] \textit{$M^TM = KI - J$;}
\item[(c)] \textit{$\hat H$ is a Hadamard matrix;}
\item[(d)] \textit{$M\mathbf{1} = 0$ and $MM^T = K\,I_D$.}
\end{itemize}

\paragraph{Proof.} (a)$\Rightarrow$(b): every column has squared
norm $D$, and equal distances between columns of equal norm mean
equal inner products, so $M^TM = \alpha I + \beta J$ with
$\alpha + \beta = D$. Distinct columns give
$\alpha = \tfrac12|v_a - v_b|^2 > 0$, so the eigenvalue $\alpha$,
of multiplicity $K - 1 = D$ on $\mathbf{1}^\perp$, makes the rank
at least $D$. The columns lie in $\mathbb{R}^D$, so the eigenvalue
$\alpha + K\beta$ on $\mathbf{1}$ vanishes. Hence $\beta = -1$ and
$\alpha = K$; this is the forced Gram matrix of the tetrad frame
with $c^2 = D$~\cite{Nelson2026frame}. (b)$\Leftrightarrow$(c):
$\hat H^T\hat H = J + M^TM$. (c)$\Leftrightarrow$(d):
$\hat H\hat H^T$ has blocks $K$, $(M\mathbf{1})^T$, $M\mathbf{1}$
and $MM^T$. (b)$\Rightarrow$(a):
$|v_a - v_b|^2 = 2K$ for $a \neq b$. \hfill $\square$

\medskip
\noindent
Every Hadamard matrix of order $K$ becomes one with first row
$\mathbf{1}^T$ after negating columns, and its remaining rows then
form a matrix $M$ whose columns are a cube simplex. Cube simplices and normalised Hadamard
matrices are therefore the same objects. Condition (d) is the pair
of conditions imposed on the tetrahedral directions in
Letter~1~\cite{Nelson2026signature}: no net direction, and no
preferred axis.

\paragraph{Lemma (Hadamard).} \textit{A Hadamard matrix of order
$n > 2$ has $4 \mid n$.}

\paragraph{Proof.} Normalise the first row to $\mathbf{1}^T$. Let
$n_{\pm\pm}$ count the columns by their signs in rows $2$ and $3$.
Orthogonality of rows $2$ and $3$ to row $1$ and to each other gives
$n_{++} + n_{+-} = n_{-+} + n_{--}$,
$n_{++} + n_{-+} = n_{+-} + n_{--}$ and
$n_{++} + n_{--} = n_{+-} + n_{-+}$, so all four counts are equal
and $n = 4\,n_{++}$. \hfill $\square$

\paragraph{Corollary 1.} \textit{A cube simplex in $Q_D$ exists
only if $D = 1$ or $4 \mid D + 1$. The first case beyond the
segment is $D = 3$.}

% ----------------------------------------------------------------------

\section*{Only Three Dimensions Fill the Cube}

\paragraph{Theorem 2.} \textit{Let $V$ be a cube simplex in $Q_D$,
$D \ge 2$. Then $V$ and $-V$ are disjoint cube simplices, and
$V \cup (-V) = Q_D$ if and only if $D = 3$. For $D = 3$, $V$ and
$-V$ are the two parity classes $\{x_1x_2x_3 = +1\}$ and
$\{x_1x_2x_3 = -1\}$, and their compound is Kepler's stella
octangula~\cite{Kepler1619,Cromwell1997}, the stellated
octahedron.}

\paragraph{Proof.} Negation preserves inner products, so $-V$ is a
cube simplex. A shared point $v_a = -v_b$ would need
$v_a \cdot v_b = -D$, but by Theorem 1 the inner products are $D$
and $-1$, and $D \ge 2$. So $|V \cup (-V)| = 2(D+1)$, which equals
$|Q_D| = 2^D$ only for $D = 3$: the equation fails at $D = 2$, and
for $D \ge 4$, $2^D > 2(D+1)$ by induction from $16 > 10$. For
$D = 3$, $v_a \cdot v_b = -1$ means that $v_a$ and $v_b$ differ in
exactly two coordinates, so they have the same product
$x_1x_2x_3$; the four points of $V$ lie in one parity class of
four, and $-V$ in the other. \hfill $\square$

\medskip
\noindent
The equation $2(D+1) = 2^D$ counts the same set twice: on the left,
the vertices of a simplex and its parity image; on the right, the
vertices of the cube.

% ----------------------------------------------------------------------

\section*{The Walsh Factorisation}

The Walsh functions on $Q_3$ are
$\chi_S(x) = \prod_{i \in S} x_i$, $S \subseteq \{1,2,3\}$, and the
Walsh matrix is $W_{S,x} = \chi_S(x)$.

\paragraph{Theorem 3.} \textit{Let $V$ be the parity class
$x_1x_2x_3 = +1$, with columns $v_0, \dots, v_3$ of $M$. Order the
points of $Q_3$ as $v_0, \dots, v_3, -v_0, \dots, -v_3$, and the
Walsh functions as $1$, $x_2x_3$, $x_1x_3$, $x_1x_2$ (even degree),
then $x_1x_2x_3$, $x_1$, $x_2$, $x_3$ (odd degree). Then}
\begin{equation}
  W = H_2 \otimes \hat H = \begin{pmatrix}
    \hat H & \hat H \\ \hat H & -\hat H \end{pmatrix},
\end{equation}
\textit{where $H_2$ has rows $(1, 1)$ and $(1, -1)$.}

\paragraph{Proof.} On $V$, $x_1x_2x_3 = 1$, so $x_2x_3 = x_1$,
$x_1x_3 = x_2$ and $x_1x_2 = x_3$: both lists restrict to
$1, x_1, x_2, x_3$, whose values on $v_0, \dots, v_3$ are the rows
of $\hat H$. Since $\chi_S(-x) = (-1)^{|S|}\chi_S(x)$, the even
functions take the same values on $-V$ and the odd functions the
opposite values. \hfill $\square$

\medskip
\noindent
The factor $H_2$ is spatial parity, $x \to -x$, acting on Walsh
functions by their degree; the factor $\hat H$ is the frame. On
each tetrahedron, every Walsh function coincides with its
complement up to a common sign, so the eight functions reduce to
four.

% ----------------------------------------------------------------------

\section*{Why Four Legs}

\paragraph{Corollary 2.} \textit{For a cube simplex with $K$ even,
each of the $K - 1$ rows of $M$ splits the $K$ points into two
halves, and the rows give $K - 1$ distinct splits. They are all such
splits if and only if $K = 2$ or $K = 4$.}

\paragraph{Proof.} By Theorem 1(d), each row has as many $+1$ as
$-1$ entries. Two rows never give the same split: they are
orthogonal, and a row and its negative are not. There are
$\tfrac12\binom{K}{K/2}$ splits, equal to $K - 1$ at $K = 2$ and
$K = 4$; for $K \ge 6$,
$\tfrac12\binom{K}{K/2} \ge \tfrac12\binom{K}{2} > K - 1$.
\hfill $\square$

\medskip
\noindent
At $K = 4$ the three rows of $M$ are the three ways of splitting
the four points, the four legs of the frame, into two pairs.

% ----------------------------------------------------------------------

\section*{Summary}

On the cube, the regular simplex of the tetrad frame and the
normalised Hadamard matrix are one object: its columns satisfy the
forced Gram matrix, and its rows satisfy closure and isotropy.
Hadamard's condition limits such simplices to $D = 1$ and
$D + 1 \equiv 0 \pmod 4$. A simplex and its parity image fill the
cube only in three dimensions, where they are the two parity
classes and form the stellated octahedron; this is what
$2(D+1) = 2^D$ counts. The Walsh matrix of the three-cube factorises
exactly as spatial parity times the frame.

% ----------------------------------------------------------------------

\begin{thebibliography}{9}

\bibitem{Nelson2026frame}
S.~Nelson,
\textit{The Spatial Geometry of the Lorentz Tetrad} (2026).

\bibitem{Nelson2026signature}
S.~Nelson,
\textit{Supplementary Material: From the Lorentz Signature to the
Tetrahedral Frame} (2026).

\bibitem{Kepler1619}
J.~Kepler,
\textit{Harmonices Mundi}
(Linz, 1619).

\bibitem{Cromwell1997}
P.~R.~Cromwell,
\textit{Polyhedra}
(Cambridge University Press, Cambridge, 1997).

\end{thebibliography}

\end{document}
